% Constraint Fields in TIDAL
% Source: docs/constraint_fields.md (migrated to LaTeX)
% Uses macros from preamble.tex

\section{Constraint Fields in TIDAL}
\label{sec:constraint-fields}

\textbf{Status:} Active \quad
\textbf{Date:} 2026-02-19 \quad
\textbf{Applies to:} Any Lagrangian that produces mixed time-derivative orders

\subsection{Overview}
\label{sec:cf-overview}

When the Wolfram pipeline derives equations of motion via \texttt{VarD} (Euler--Lagrange
variation), the resulting component equations may have different time-derivative
orders on the LHS:

\begin{table}[htbp]
\centering
\caption{Classification of component equations by time-derivative order.}
\label{tab:cf-classification}
\begin{tabular}{@{}clp{4cm}p{4.5cm}@{}}
\toprule
\texttt{time\_derivative\_order} & Type & Mathematical character & Solver treatment \\
\midrule
2 & \textbf{Dynamical} (wave) & Hyperbolic: $\ddt^2 \varphi = \mathrm{RHS}$ & ODE-integrated (CVODE / IDA / scipy) \\
1 & \textbf{First-order} & Parabolic/advective: $\ddt \varphi = \mathrm{RHS}$ & ODE-integrated \\
0 & \textbf{Constraint} & Elliptic: $0 = \mathrm{RHS}$ & Algebraic in IDA; pre-solved (FFT / sparse) \\
\bottomrule
\end{tabular}
\end{table}

Constraint fields are \textbf{not} integrated forward in time. They are determined
by an elliptic equation at each instant, given the current state of the
dynamical fields. This is a fundamental feature of the Lagrangian, not a
pipeline limitation.

\subsection{Physics: Why Constraint Equations Arise}
\label{sec:cf-physics}

\subsubsection{The Proca Example}
\label{sec:cf-proca}

Consider the Proca Lagrangian for a massive vector field $A_\mu$ in 2+1D with
Minkowski signature $(-,+,+)$:

\begin{equation}
  \lag = -\tfrac{1}{2}\,(\partial_a A_b)\,\eta^{ac}\,\eta^{bd}\,(\partial_c A_d)
       + \tfrac{1}{2}\,(\partial_a A_b)\,\eta^{ac}\,\eta^{bd}\,(\partial_d A_c)
       - \tfrac{m^2}{2}\,A_a\,\eta^{ab}\,A_b
  \label{eq:proca-lagrangian}
\end{equation}

The Euler--Lagrange equation $\partial_b F^{ba} + m^2 A^a = 0$ yields three component
equations. For $a = 0$ (the temporal component):

\begin{equation}
  \ddx F^{x0} + \ddy F^{y0} + m^2 A^0 = 0
  \label{eq:proca-temporal}
\end{equation}

Expanding the field strength $F^{i0} = \eta^{ii}\,\eta^{00}\,(\partial_i A_0 - \ddt A_i)$:

\begin{equation}
  -\ddx(\ddx A_0 - \ddt A_1) - \ddy(\ddy A_0 - \ddt A_2) - m^2 A_0 = 0
  \label{eq:proca-expanded}
\end{equation}

Rearranging:

\begin{equation}
  \lapl A_0 - m^2 A_0 = -\ddx(\ddt A_1) - \ddy(\ddt A_2)
  \label{eq:proca-helmholtz}
\end{equation}

\textbf{There is no $\ddt^2 A_0$ term.} The metric signature $\eta^{00} = -1$ causes
the temporal kinetic term to vanish from the $a = 0$ equation. The result is a
\textbf{Helmholtz equation} --- an elliptic PDE that determines $A_0$ at each instant
from the momenta of the spatial components $A_1$, $A_2$.

\subsubsection{Analogy: Coulomb's Equation in Electrostatics}
\label{sec:cf-coulomb}

This is analogous to the electric potential $V$ in electrostatics. Poisson's
equation $\lapl V = -\rho/\varepsilon_0$ determines $V$ instantaneously from the charge density
$\rho$. If $\rho$ changes (charges move), $V$ changes --- but $V$ has no independent
dynamics. It is a ``slave'' to the dynamical degrees of freedom.

The same pattern appears in:

\begin{itemize}
  \item \textbf{Numerical relativity:} The Hamiltonian and momentum constraints determine
    the lapse and shift (Baumgarte \& Shapiro, \emph{Numerical Relativity}, Cambridge,
    2010, Chapter~3).
  \item \textbf{Incompressible fluid dynamics:} The pressure $p$ satisfies a Poisson
    equation $\lapl p = -\rho\,\partial_i u_j\,\partial_j u_i$ and is determined at each instant from
    the velocity field (Chorin, 1968; Temam, 1969).
  \item \textbf{Gauge field theories:} The temporal component $A_0$ of any gauge field
    satisfies a constraint (Gauss's law), not a wave equation.
\end{itemize}

\subsubsection{General Rule}
\label{sec:cf-general-rule}

For any Lagrangian with Lorentzian signature, the EOM for the \textbf{temporal
component} of a vector or tensor field will typically have
\texttt{time\_derivative\_order = 0}. This is not an error --- it is a direct consequence
of the signature. The pipeline correctly identifies and classifies these
equations.

\subsection{Solver Architecture}
\label{sec:cf-solver}

\subsubsection{State Vector Layout}
\label{sec:cf-state-layout}

Constraint fields are included in the ODE state vector so that dynamical
equations can reference them. The state layout for coupled Proca
($A_0, A_1, \pi_1, A_2, \pi_2, B_0, B_1, \pi_4, B_2, \pi_5$) includes both
constraint fields ($A_0$, $B_0$) and dynamical fields with their momenta.

\subsubsection{Solver Treatment}
\label{sec:cf-solver-treatment}

\paragraph{With IDA (DAE solver --- default for systems with constraints)}
\label{sec:cf-ida}

Constraint equations are treated as algebraic equations within the DAE
system. IDA's Newton iteration simultaneously solves for all field values
(dynamical + constraint) at each timestep. The constraint pre-solve module
(\texttt{tidal/solver/constraint\_solve.py}) ensures consistent initial conditions
before IDA starts. See \texttt{tidal/solver/ida.py} for the residual construction.

\textbf{Residual structure:} For constraint fields, the IDA residual is:
\begin{lstlisting}
F_i = RHS_i(state)   (algebraic: no time derivative on LHS)
\end{lstlisting}
For dynamical fields:
\begin{lstlisting}
F_i = yp_i - RHS_i(state)   (differential: yp = dy/dt)
\end{lstlisting}

\paragraph{Unified constraint IC solver}
\label{sec:cf-unified-ic}

\texttt{ensure\_consistent\_ic()} in \texttt{tidal/solver/constraint\_solve.py} is a single entry point that handles
all constraint equation types --- standard, subsidiary, and self-term --- in
a unified iterative framework:

\begin{enumerate}
  \item \textbf{Phase 1 --- Standard constraints} (\texttt{constraint\_solver.enabled=True}):
    Solves for the constraint field using the three-tier solver (FFT $\to$ sparse
    matrix $\to$ auto-select). Handles coupled constraints (e.g., coupled Proca
    $A_0 \leftrightarrow B_0$). This is the same behavior as the previous \texttt{pre\_solve\_constraints()}.

  \item \textbf{Phase 2 --- Iterative propagation} (remaining constraints):
    For each constraint equation $0 = \sum c_i \cdot \mathrm{op}_i(\text{field}_i)$, identifies
    which fields are ``determined'' (non-zero IC or already solved) vs ``free''
    (zero). When exactly one free field remains, solves for it. Iterates until
    no more progress (cascade resolution).

  \item \textbf{Phase 3 --- Verification}: Evaluates all constraint residuals. In strict
    mode (default), raises \texttt{ValueError} if any constraint is violated. With
    \texttt{--allow-inconsistent-ic}, issues a warning instead.
\end{enumerate}

\textbf{Solver tiers} (used in both phases 1 and 2):
\begin{enumerate}
  \item \textbf{Tier 1 (FFT)} --- periodic BCs, constant coefficients ($O(N \log N)$)
  \item \textbf{Tier 2 (Sparse probe)} --- non-periodic BCs or variable coefficients ($O(N^2)$ build, $O(N)$ solve)
  \item \textbf{Automatic selection} --- \texttt{\_select\_method()} chooses the fastest applicable tier
\end{enumerate}

\textbf{Gauge regularisation} for singular Poisson with periodic BCs: the zero
Fourier mode is pinned ($\hat{u}_{0,\ldots,0} = 0$) to fix the gauge freedom.
This is numerical regularisation, not physics gauge fixing --- observables
depend on derivatives of the constraint field, not its absolute value.

\subsubsection{Subsidiary Constraint Propagation}
\label{sec:cf-subsidiary}

Some constraint equations involve dynamical fields but not the constraint
field itself (no self-term). These arise from gauge conditions:

\begin{table}[htbp]
\centering
\caption{Constraint equation types.}
\label{tab:cf-constraint-types}
\begin{tabular}{@{}llp{6cm}@{}}
\toprule
Constraint type & Example & Equation form \\
\midrule
Standard (self-term) & EM $A_0$ & \texttt{laplacian($A_0$) + source = 0} \\
Subsidiary (no self-term) & GW transverse & \texttt{gradient\_x($h_4$) + gradient\_x($h_7$) = 0} \\
Self-term algebraic & GW traceless & \texttt{identity($h_4$) + identity($h_7$) + identity($h_9$) = 0} \\
\bottomrule
\end{tabular}
\end{table}

The unified solver handles all three cases:

\begin{itemize}
  \item \textbf{Standard}: target = constraint field (e.g., $A_0$), solve $L[A_0] = -\text{source}$
  \item \textbf{Subsidiary}: target = the one free dynamical field (e.g., $h_7$),
    solve $L[h_7] = -\text{source}$ where source comes from determined fields
  \item \textbf{Self-term algebraic}: target = constraint field with self-terms (e.g., $h_9$),
    solve $\mathrm{identity}(h_9) = -(h_4 + h_7)$
\end{itemize}

Cascade example (GW TT gauge, \texttt{--ic-component h\_4}):

\begin{enumerate}
  \item User sets $h_4 = \text{Gaussian}$
  \item Transverse constraint determines $h_7 = -\text{Gaussian}$
  \item Traceless constraint determines $h_9 = 0$
  \item Velocity constraints auto-satisfied (all velocities start at zero)
\end{enumerate}

\subsection{Constraint Momentum Estimation}
\label{sec:cf-momentum}

\subsubsection{The Problem}
\label{sec:cf-momentum-problem}

Although constraint fields have no evolution equation, they \textbf{do change in
time} because their source terms (dynamical field momenta) evolve. If a
dynamical equation references $\ddx(\pi_0)$ --- i.e., $\ddx(\ddt A_0)$ ---
this term should be nonzero.

Prior to this fix, the code set \texttt{virtual\_momenta["A\_0"] = 0} unconditionally,
dropping the $\ddx(\ddt A_0)$ coupling from the dynamical equations.

\subsubsection{The Fix: Analytical Constraint Momentum via Time Differentiation}
\label{sec:cf-momentum-fix}

Rather than estimating $\pi_N$ via finite differences (which introduces a time
lag that causes $k$-dependent instability through the constraint-dynamical
feedback loop), we compute $\pi_N$ \textbf{exactly} by differentiating the constraint
equation in time.

\paragraph{Mathematical Derivation}
\label{sec:cf-math-derivation}

Given a constraint equation of the form:
\begin{equation}
  L[A_0] = -S(\pi_1, \pi_2, \ldots, B_0, \ldots)
  \label{eq:cf-constraint-form}
\end{equation}
where $L$ is the self-operator (e.g., $\lapl - m^2$) and $S$ is the source
(terms referencing other fields and their momenta), differentiating both sides
with respect to time yields:
\begin{equation}
  L[\ddt A_0] = L[\pi_0] = -\ddt S
  \label{eq:cf-differentiated}
\end{equation}

This is another elliptic equation for $\pi_0$, with the \textbf{same self-operator}
$L$ as the original constraint, and a new source $-\ddt S$ computed from the
time derivatives of the dynamical fields.

\paragraph{Computing $\ddt S$}
\label{sec:cf-dtS}

The time derivative of each source term is obtained by replacing fields with
their time derivatives:

\begin{table}[htbp]
\centering
\caption{Time-derivative replacement rules for source terms.}
\label{tab:cf-replacements}
\begin{tabular}{@{}lll@{}}
\toprule
Source term type & Replacement & Source of replacement \\
\midrule
$c \cdot \mathrm{op}(\pi_j)$ where $j$ is 2nd-order & $c \cdot \mathrm{op}(\ddt \pi_j) = c \cdot \mathrm{op}(\mathrm{RHS}_j)$ & Dynamical equation $j$ \\
$c \cdot \mathrm{op}(h_j)$ where $h_j$ is 2nd-order field & $c \cdot \mathrm{op}(\pi_j)$ & Momentum from state \\
$c \cdot \mathrm{op}(\text{constraint\_field}_j)$ & $c \cdot \mathrm{op}(\pi_j)$ & Off-diagonal coupling \\
\bottomrule
\end{tabular}
\end{table}

\paragraph{Handling Back-Coupling}
\label{sec:cf-back-coupling}

The dynamical equation $\mathrm{RHS}_j$ may itself contain $\nabla(\pi_0)$ terms
(the very momentum we are solving for). These terms are moved to the LHS,
creating an \textbf{augmented operator}:

\begin{equation}
  L_\mathrm{aug}[\pi_0] = L[\pi_0] + \text{back\_coupling}[\pi_0] = -\ddt S_\mathrm{base}
  \label{eq:cf-augmented}
\end{equation}
where $\ddt S_\mathrm{base}$ is computed with all constraint momenta set to zero.

The back-coupling contribution is a product of Fourier multipliers. For
example, in coupled Proca:

\begin{itemize}
  \item Source term: $-1 \cdot \ddx(\pi_1)$ in $A_0$ constraint
  \item $A_1$ dynamical eq contains: $-1 \cdot \ddx(\pi_0)$
  \item Back-coupling multiplier: $(-1) \cdot ik_x \cdot (-1) \cdot ik_x = -k_x^2$
  \item Similarly for $y$: $(-1) \cdot ik_y \cdot (-1) \cdot ik_y = -k_y^2$
  \item Total: augmented diagonal $= (-k^2 - m^2) + (-k^2) = -2k^2 - m^2$
\end{itemize}

\paragraph{Block FFT Solve for Coupled Constraints}
\label{sec:cf-block-fft}

When multiple constraint fields are coupled (e.g., $A_0 \leftrightarrow B_0$ in coupled
Proca), the system is solved simultaneously via batched SVD:

\begin{equation}
  M_\mathrm{aug}(\bm{k}) \cdot [\pi_{A_0}, \pi_{B_0}, \ldots]^T = -[\ddt S_{A_0}, \ddt S_{B_0}, \ldots]^T
  \label{eq:cf-block-system}
\end{equation}
where $M_\mathrm{aug}(\bm{k})$ has:
\begin{itemize}
  \item \textbf{Diagonal entries:} self-operator multiplier + back-coupling
  \item \textbf{Off-diagonal entries:} from cross-constraint coupling (e.g., \texttt{gcoup})
\end{itemize}

This uses the same SVD infrastructure and Tikhonov regularization as the
existing block FFT constraint solver (\texttt{\_solve\_coupled\_constraints\_fft}).

\subsubsection{Why This Is Better Than Finite Differences}
\label{sec:cf-why-analytical}

The previous approach estimated $\pi_N$ via backward FD:
$\pi_N \approx (\text{field}(t) - \text{field}(t_\mathrm{prev})) / \Delta t$. This introduced a \textbf{time lag} that
turned the constraint-dynamical feedback loop into a delay differential
equation (DDE). Perturbative analysis of the DDE characteristic equation
showed a growth rate $\sigma \propto k^4 \Delta t / m^2$ --- exponential blow-up of high-$k$ modes
that no time step reduction could cure.

The analytical approach eliminates the time lag entirely:

\begin{itemize}
  \item \textbf{No persistent state:} The method is a pure function of $(\text{state}, t)$,
    inherently compatible with adaptive ODE solvers (no anchor/caching needed).
  \item \textbf{Exact at every instant:} No $O(\Delta t)$ error --- the momentum is the exact
    solution of the differentiated constraint equation.
  \item \textbf{$k$-independent stability:} The augmented operator $L_\mathrm{aug}$ has strictly
    negative eigenvalues for all $k$, preventing instability at any wavelength.
\end{itemize}

\subsubsection{Implementation}
\label{sec:cf-implementation}

Method \texttt{\_compute\_analytical\_constraint\_momenta()} runs after the constraint
solve, computing virtual momenta via \texttt{\_solve\_constraint\_momenta\_fft()}.

The method is theory-agnostic: it reads the equation structure from the JSON
spec and automatically identifies source terms, back-coupling chains, and
cross-constraint coupling without any theory-specific knowledge.

\subsubsection{Limitations}
\label{sec:cf-limitations}

\begin{itemize}
  \item \textbf{Periodic BCs only:} Requires FFT. For non-periodic grids, a future
    matrix-based analytical solve extension is needed.
  \item \textbf{Position-dependent self-term coefficients:} Same limitation as the FFT
    constraint solver --- raises \texttt{ValueError} suggesting \texttt{method='matrix'}.
    Falls back to $\pi = 0$ with a debug-level log message.
  \item \textbf{Time-dependent source coefficients:} $\ddt(c(t) \cdot \mathrm{op}(\text{field}))$ requires
    $\ddt c$, which is not yet implemented. Raises \texttt{ValueError} if detected.
    (No current example has this case.)
\end{itemize}

\subsection{References}
\label{sec:cf-references}

\begin{enumerate}
  \item Baumgarte, T.\ W.\ \& Shapiro, S.\ L.\ (2010). \emph{Numerical Relativity: Solving
    Einstein's Equations on the Computer}. Cambridge University Press.
    --- Chapters 2--3: constraint equations in 3+1 decomposition.

  \item Chorin, A.\ J.\ (1968). ``Numerical solution of the Navier--Stokes equations.''
    \emph{Mathematics of Computation}, 22(104), 745--762.
    --- Projection method: pressure as an elliptic constraint.

  \item Jackson, J.\ D.\ (1999). \emph{Classical Electrodynamics} (3rd ed.). Wiley.
    --- Section 6.3: Coulomb gauge and the scalar potential as a constraint.

  \item Ruehl, W.\ \& Yurov, V.\ P.\ (2004). ``Proca equations derived from first
    principles.'' arXiv:hep-th/0412059.
    --- Structure of Proca EOM and constraint analysis.
\end{enumerate}
